% =====================================================================
\chapter{Bevezetés}
\label{ch:bevezetes}
% =====================================================================

\section{Motiváció}

A felsőoktatásban és a középiskolában a feleletválasztós kérdés (\emph{multiple-choice question}, MCQ)
a tudásmérés egyik legszélesebb körben használt eszköze: gépi úton javítható, nagy létszámú kurzusokon is
skálázható, és jól illeszkedik a tanulásmenedzsment-rendszerek (pl.\ Moodle) kvízmoduljaihoz. Előállítása
azonban drága. A pszichometriai szakirodalom szerint egy jó kérdés törzse önmagában is érthető és
egyértelmű, pontosan egy válaszlehetősége helyes a tananyag szerint, a disztraktorok pedig azonos
témájúak, formailag hasonlóak és egy felkészületlen hallgató számára hihetők, ám igazolhatóan
hibásak~\cite{haladyna2002review,gierl2017distractors}. Egy tapasztalt oktató egy-egy ilyen kérdésen
perceket, egy teljes vizsgán órákat dolgozik, és a disztraktorok megírása a munka legnehezebb része.

A nagy nyelvi modellek (\emph{large language model}, LLM) megjelenése~\cite{brown2020gpt3} gyökeresen
átalakította ezt a munkafolyamatot. Az oktatástechnológiai irodalom egyszerre hangsúlyozza a lehetőségeket
és a kockázatokat~\cite{kasneci2023chatgpt}: a modellek képesek kérdéseket és disztraktorokat
fogalmazni~\cite{kurdi2020systematic,alhazmi2024distractor}, ugyanakkor hajlamosak folyékony, de a forrás
által nem alátámasztott állításokat generálni~\cite{ji2023hallucination}. A gyakorlatban a
legelterjedtebb használati mód a legegyszerűbb: az oktató feltölti az előadásjegyzetét egy felhőben futó
csevegő asszisztensbe, és „tíz feleletválasztós kérdést” kér. Ezt a munkafolyamatot a továbbiakban
\emph{status quo}-nak vagy \emph{közvetlen megközelítésnek} nevezzük.

\section{Problémafelvetés}

A közvetlen megközelítés kényelmes, de az iskolai alkalmazás szempontjából négy, egymástól független
gyengesége van. Ezek adják dolgozatunk kiindulópontját.

\begin{description}[leftmargin=1.1cm,labelwidth=0.9cm,style=sameline,format=\normalfont\bfseries]
\item[(P1)] \textbf{A helyesség nincs ellenőrizve.} Semmi sem garantálja, hogy a megjelölt kulcsot a
dokumentum alátámasztja, és hogy a disztraktorok valóban hamisak. Az LLM-ek hihető, de alá nem támasztott
állításokat állítanak elő~\cite{ji2023hallucination}; az önellenőrzés ezt csökkentheti, de nem
szünteti meg~\cite{madaan2023selfrefine,dhuliawala2024cove,huang2024selfcorrect}.
\item[(P2)] \textbf{A szerkezet és a lefedettség nincs kontrollálva.} A kért darabszám, típus és nehézség
csak kérés a promptban. A hosszú bemenetet a modellek egyenetlenül használják: a kontextus elejére és
végére jobban figyelnek, mint a közepére~\cite{liu2024lost}, a névleges kontextushossznál jóval rövidebb
az a hossz, amelyen még megbízhatóan dolgoznak~\cite{hsieh2024ruler}. Ennek következménye, hogy a kérdések
a dokumentum néhány szakaszára tömörülnek.
\item[(P3)] \textbf{Adatvédelem.} A tananyag személyes adatot (pl.\ esettanulmányban szereplő neveket,
hallgatói munkákat) is tartalmazhat; a felhőszolgáltatás az adatkezelés címzettjévé
válik~\cite{gdpr2016}, az intézmény pedig elveszíti a kontrollt a tartalom sorsa felett.
\item[(P4)] \textbf{Függés a nagy, távoli modellektől.} A teljes dokumentum egyetlen promptban hosszú
kontextust és nagy modellt igényel. Az iskola által saját hardveren futtatható kis modellek (néhány
milliárd paraméter, 8\,GB videomemória) kontextusa és memóriája korlátos.
\end{description}

Fontos látni, hogy a négy probléma nem oldható meg egyszerűen „jobb prompttal”. A P1 és a P2 a
\emph{módszer} hiányából fakad: ha a vizsga szerkezete, a kérdések forrása és az ellenőrzés nem explicit
lépések, akkor nem is mérhetők és nem is kényszeríthetők ki. A P3 és a P4 pedig egymást erősítik: a
helyi futtatás megoldaná az adatvédelmi problémát, de a helyben futtatható kis modellek éppen a hosszú
kontextusú, egylépéses feladatban a leggyengébbek. Kutatásunk alapgondolata ezért az, hogy a
feladatot rövid, jól specifikált, egyenként ellenőrizhető lépésekre bontva egy kis, helyben futó modell
is képes lehet megbízható vizsgát előállítani, és ez a felbontás egyben a P1 és P2 problémát is kezeli.

\section{Kutatási kérdések és hipotézisek}

A fenti problémák alapján három kutatási kérdést fogalmaztunk meg.

\begin{description}[leftmargin=1.3cm,labelwidth=1.1cm,style=sameline,format=\normalfont\bfseries]
\item[RQ1] A strukturált generálás -- tervezés, kérdésenkénti forráshoz kötés és ellenőrzés -- érvényesebb
és jobban lefedő vizsgát eredményez-e, mint ha a modellnek a dokumentumot (vagy annak visszakeresett
részleteit) egyetlen promptban adjuk át?
\item[RQ2] Mennyivel járul hozzá a minőséghez az egyes komponens (tervező, ellenőrző, fogalomgráf,
hibrid visszakeresés, darabolási stratégia), és mennyibe kerül futási időben, modellhívásban és
tokenben?
\item[RQ3] Mennyire közelíti meg egy 8\,GB-os fogyasztói GPU-n futó, 4 bitre kvantált 7B modell egy
nagy, szerveren futó modell teljesítményét?
\end{description}

A kutatási kérdésekhez a következő hipotéziseket rendeljük. Értékelésüket a pilot és a fő vizsgálat
alapján \az{\ref{sec:hipotezis-ertekeles}}.~szakasz adja.

\begin{hipotezis}[Tervezés]\label{h:terv}
Azonos generátormodell mellett a vizsgaterv (\arm{E1}) a naiv RAG alapvonalhoz (\arm{E0}) képest
növeli a kérdések alátámasztottságát és a vizsga darablefedettségét, a formai megfelelés romlása nélkül.
\end{hipotezis}

\begin{hipotezis}[Ellenőrzés]\label{h:verif}
Az ellenőrző lánc (\arm{E2}) a tervezéshez (\arm{E1}) képest tovább növeli az alátámasztottságot és a
vak megoldhatóságot; az ellenőrző ítélete pozitívan egyezik a független bíró ítéletével ($\kappa>0$).
\end{hipotezis}

\begin{hipotezis}[Lefedettség]\label{h:lefed}
A teljes dokumentumot egyetlen promptban kapó alapvonal (\arm{B-doc}) lefedettsége -- modellmérettől
függetlenül -- alacsonyabb, mint a tervezett változatoké.
\end{hipotezis}

\begin{hipotezis}[Helyi vs.\ szerver]\label{h:helyi}
A helyi 7B modell a tervezett és ellenőrzött csővezetékben a szervermodell naiv változatának
minőségi indexét megközelíti (különbség $<0{,}1$), miközben a GPU-memóriaigénye 8\,GB alatt marad.
\end{hipotezis}

\begin{hipotezis}[Mérés]\label{h:meres}
A generátoroktól eltérő modellcsaládba tartozó LLM-bíró ítéletei az oktatói értékelésekkel legalább
elfogadható megbízhatósággal egyeznek (Krippendorff-$\alpha\ge 0{,}667$).
\end{hipotezis}

\section{Hozzájárulások}

Dolgozatunk négy, egymásra épülő hozzájárulást tartalmaz.

\begin{enumerate}[label=(\arabic*)]
\item \textbf{A \Mimir{} generálási módszer} (\ref{ch:modszer}.~fejezet): ablakozott mondatbeágyazáson
alapuló szemantikus darabolás, lefedettségvezérelt vizsgaterv a legnagyobb maradék elvén kiosztott
Bloom-szintekkel, slotonkénti sűrű vagy hibrid visszakeresés, forrásrészletekre hivatkozó egyenkénti
kérdésgenerálás, legolcsóbbtól a legdrágábbig rendezett ellenőrző lánc visszacsatolással és
tartalékfogalmakkal, valamint egy munkamenet-hatókörű fogalomgráf (GraphRAG-lite). A módszert formálisan
is megfogalmazzuk (\ref{ch:formalizmus}.~fejezet), és algoritmikusan leírjuk.
\item \textbf{Adatvédelmet biztosító helyi architektúra} (\ref{ch:architektura}.~fejezet): lazán
csatolt mikroszolgáltatások, PostgreSQL-alapú zárolásmentes feladatsor, zéró adatmegőrzés és
ujjlenyomat-alapú audit, valamint a GPU-memóriaigény analitikus modellje, amelyet mérésekkel
validálunk.
\item \textbf{Kiértékelési módszertan és nyílt keretrendszer} (\ref{ch:kiertekeles}.~fejezet): tizenöt
egytényezős kísérleti kar, kétnyelvű adatkészlet referenciakérdésekkel, formálisan definiált mutatók,
modellcsaládon kívüli LLM-bíró, oktatói értékelési protokoll és dokumentumszintű statisztika.
\item \textbf{Empirikus eredmények} (\ref{ch:eredmenyek}.~fejezet): egy pilot vizsgálat, amely két
korábbi méréseket érvénytelenítő mérési hibát, a válaszszivárgást mint új hibamódot és a futásról futásra
jelentkező zajt tárta fel, valamint egy 14 karos, kilenc dokumentumos fő vizsgálat, amely szerint a helyi
7B modellel futó tervezett és ellenőrzött csővezeték a közös dokumentumokon a naiv RAG-ot ($0{,}64$) és a
teljes dokumentumos promptot ($0{,}83$) is felülmúlta ($0{,}96$), és elérte a teljes dokumentumot kapó
szervermodellt ($0{,}95$).
\end{enumerate}

\section{A dolgozat felépítése}

\Az{\ref{ch:szakirodalom}}.~fejezet áttekinti az automatikus kérdésgenerálás, a visszakereséssel
kiegészített generálás, az ellenőrzés és a nyelvimodell-alapú kiértékelés szakirodalmát, és
azonosítja a kutatási rést. \Az{\ref{ch:formalizmus}}.~fejezet a feladatot formálisan definiálja,
bevezeti a minőségi predikátumokat és a célfüggvényt. \Az{\ref{ch:architektura}}.~fejezet a rendszer
architektúráját, adatvédelmi mechanizmusait és erőforrás-modelljét mutatja be. \Az{\ref{ch:modszer}}.~fejezet a generálási módszer algoritmusait és azok költségelemzését tárgyalja.
\Az{\ref{ch:kiertekeles}}.~fejezet a kiértékelési módszertant, \az{\ref{ch:eredmenyek}}.~fejezet a pilot és a fő
vizsgálat eredményeit ismerteti. \Az{\ref{ch:diszkusszio}}.~fejezet az eredmények értelmezését és az érvényességet
veszélyeztető tényezőket, \az{\ref{ch:osszefoglalas}}.~fejezet az összefoglalást és a további munkát
tartalmazza. A függelék a promptokat, a konfigurációkat és a reprodukálás lépéseit közli.
